\problema{Jump Game}{55}{src/07-dp/55-jump-game.cpp}{M}

Dado un arreglo \texttt{nums} donde cada elemento \texttt{nums[i]}
representa el número máximo de posiciones que se puede saltar hacia
adelante desde el índice \texttt{i}, determina si es posible llegar desde el
índice \texttt{0} hasta el último índice del arreglo.

\paragraph{La idea, con un ejemplo de la vida real.} Imagina piedras
colocadas en fila para cruzar un río, y cada piedra tiene escrito cuántas
piedras hacia adelante puedes brincar desde ahí como máximo. En vez de
probar todas las secuencias posibles de saltos (que se ramifican
exponencialmente), basta con llevar la cuenta de \emph{cuál es la piedra más
lejana a la que se puede llegar}, considerando todo lo que ya se sabe
alcanzable. Si en algún momento la siguiente piedra que toca revisar está
más allá de ese alcance máximo, ya no hay forma de llegar ahí y el cruce es
imposible.

\begin{center}
\ttfamily
\begin{tabular}{l}
nums = [2, 3, 1, 1, 4]\\
desde el índice 0 se alcanza el índice 2\\
desde el índice 1 se alcanza el índice 4 (el último) -> true\\
\end{tabular}
\end{center}

\paragraph{Cómo lo resuelve el código, paso a paso.} Este problema admite
una solución de programación dinámica (\emph{dynamic programming}, DP)
tabular, pero la solución esperada en entrevista es un algoritmo
\textbf{greedy} (voraz), más simple y con espacio constante.
\begin{enumerate}
  \item Mantenemos una sola variable, \texttt{alcanceMaximo}, inicializada
        en \texttt{0}.
  \item Recorremos el arreglo de izquierda a derecha con un índice
        \texttt{i}.
  \item Si \texttt{i} es mayor que \texttt{alcanceMaximo}, significa que
        ninguna posición ya confirmada como alcanzable puede llegar hasta
        aquí: devolvemos \texttt{false} de inmediato.
  \item Si no, actualizamos \texttt{alcanceMaximo = max(alcanceMaximo, i +
        nums[i])}, es decir, consideramos si saltar desde la posición actual
        extiende el alcance máximo conocido.
  \item Si el ciclo termina sin encontrar una posición inalcanzable,
        devolvemos \texttt{true}.
\end{enumerate}

\lstinputlisting{src/07-dp/55-jump-game.cpp}

\complejidad{$O(n)$}{$O(1)$}

\paragraph{¿Qué significa esa complejidad?} Recorremos el arreglo una sola
vez, actualizando una sola variable en cada paso, así que el tiempo es
lineal en el tamaño del arreglo. El espacio es constante porque no se
necesita ninguna tabla ni estructura auxiliar, a diferencia de una solución
DP tabular que gastaría espacio lineal.

\paragraph{Consejos para la entrevista (Amazon).} Vale la pena mencionar
primero la versión DP (\texttt{dp[i]} = ``¿se puede llegar al índice
\texttt{i}?'', mirando hacia atrás todas las posiciones que podrían saltar
hasta ahí) para demostrar que se entiende el problema desde ese ángulo, y
luego señalar la mejora greedy y por qué es válida: \texttt{alcanceMaximo}
resume toda la información relevante de las posiciones ya visitadas, sin
necesidad de recordar el detalle de cada una. Un caso límite importante es
un cero en medio del arreglo que \emph{no} bloquea el camino porque una
posición anterior ya tenía alcance suficiente para saltar por encima de él;
hay que distinguirlo del caso en que el cero sí bloquea porque ninguna
posición previa alcanza a rebasarlo.
